\chapter{Conclusiones}\label{cap:conclusiones}

El desarrollo integral de este sistema logístico ha servido para validar de forma empírica la idoneidad de las bases de datos de grafos para la resolución de problemas de enrutamiento que, históricamente, penalizan severamente el rendimiento en sistemas de naturaleza relacional o tabular.

El uso de un diseño multimodelo en ArangoDB ha simplificado drásticamente el modelo mental del dominio de negocio. Paralelamente, la adopción del Patrón Repositorio en la capa de software ha asegurado un código legible y escalable. La contenedorización completa del ciclo de vida del proyecto mediante Docker me ha permitido ofrecer una solución ``llave en mano'', mitigando la clásica problemática de las dependencias locales y garantizando un despliegue transparente para su evaluación.

Aunque la curva de aprendizaje inicial del lenguaje AQL y del razonamiento transversal en grafos fue pronunciada, los resultados técnicos conseguidos demuestran una capacidad resolutiva sólida y justifican plenamente la elección tecnológica, logrando una plataforma altamente tolerante a fallos logísticos.

\section*{Uso de Herramientas de IA}
En cumplimiento con las directrices de la asignatura "Complementos de Bases de Datos", declaramos el uso de modelos generativos de IA (Gemini, de Google) como herramienta de apoyo técnico durante el desarrollo del proyecto. Sus usos específicos han sido:
\begin{itemize}
    \item \textbf{Análisis Arquitectónico y Refactorización:} Soporte consultivo en la transición del código de un modelo centralizado a uno distribuido (Patrón Repositorio) para la separación de responsabilidades en la base de datos.
    \item \textbf{Depuración de Despliegue:} Asistencia técnica puntual en la comprensión de jerarquías de directorios dentro de Docker para resolver el fallo de importación circular (\texttt{PYTHONPATH}).
    \item \textbf{Formato y Estilo:} Apoyo en la adaptación y formateo de la documentación técnica al lenguaje de marcado LaTeX.
\end{itemize}

Como autores de la entrega, declaramos que todo el código fuente y el contenido teórico generado ha sido analizado de manera crítica, depurado y validado empíricamente a través de nuestras pruebas. Comprendoemos en profundidad cada bloque lógico y componente del sistema expuesto en esta memoria.